% ==================== 附录：完整可运行源程序 ====================
% 按 format2026 第五条要求，附录包含建模所用到的完整、可运行的源程序。
% 本文件由 scripts/gen_appendix_source.py 自动生成，勿手改。
% 源码直接取自支撑材料文件夹 final_version/，故与所提交文件逐字一致。
\section{建模源程序（完整可运行）}
\label{sec:appendix-source}

以下按问题列出\textbf{生成论文全部数值结果}所需源程序。
这些代码\textbf{直接取自支撑材料文件夹 \texttt{final\_version/}}，
因此与所提交文件\textbf{逐字一致}；
运行环境与复现命令见 \S\ref{sec:appendix-env}。
全部程序使用 Python 语言与 \texttt{numpy}/\texttt{scipy}（HiGHS 求解器），
Excel 读写使用 \texttt{openpyxl}；未使用 SPSS 等需手工交互的软件，
故无交互命令需要单独记录。

\textbf{收录边界}：本文的图表由配图脚本从上述程序保存的结果文件
（\texttt{*.npz}/\texttt{*.csv}）重建，属于表达层，
不参与任何数值计算，也不影响任何结果；
按“仅收录与建模结果直接相关的源程序”的原则，
\texttt{plot\_*.py} 与 \texttt{common/plotting/} 未列入本附录，
其完整文件已随支撑材料一并提交。

\begingroup
\footnotesize
\setlength{\parindent}{0pt}
\setlength{\parskip}{0pt}

\subsection{问题一：确定性调度}

\subsubsection*{q1\_model.py}
\verbatiminput{../final_version/Q1/q1_code/q1_model.py}

\subsubsection*{run\_q1.py}
\verbatiminput{../final_version/Q1/q1_code/run_q1.py}

\subsection{问题二：滚动随机日前调度}

\subsubsection*{rolling\_window.py}
\verbatiminput{../final_version/Q2/code/rolling_window.py}

\subsubsection*{run\_experiment.py}
\verbatiminput{../final_version/Q2/code/run_experiment.py}

\subsubsection*{verify\_window.py}
\verbatiminput{../final_version/Q2/code/verify_window.py}

\subsubsection*{check\_information.py}
\verbatiminput{../final_version/Q2/code/check_information.py}

\subsubsection*{report\_experiment.py}
\verbatiminput{../final_version/Q2/code/report_experiment.py}

\subsubsection*{make\_target\_tables.py}
\verbatiminput{../final_version/Q2/code/make_target_tables.py}

\subsection{问题三：预报融合与场景树多时点调整}

\subsubsection*{q3\_data.py}
\verbatiminput{../final_version/Q3/code/q3_data.py}

\subsubsection*{q3\_forecast.py}
\verbatiminput{../final_version/Q3/code/q3_forecast.py}

\subsubsection*{q3\_model.py}
\verbatiminput{../final_version/Q3/code/q3_model.py}

\subsubsection*{run\_q3.py}
\verbatiminput{../final_version/Q3/code/run_q3.py}

\subsubsection*{verify\_q3.py}
\verbatiminput{../final_version/Q3/code/verify_q3.py}

\subsection{问题四-2：波动电价下的日前调度}

\subsubsection*{q42\_model.py}
\verbatiminput{../final_version/Q4/part2/code/q42_model.py}

\subsubsection*{run\_q42.py}
\verbatiminput{../final_version/Q4/part2/code/run_q42.py}

\subsubsection*{verify\_q42.py}
\verbatiminput{../final_version/Q4/part2/code/verify_q42.py}

\subsubsection*{check\_information.py}
\verbatiminput{../final_version/Q4/part2/code/check_information.py}

\subsubsection*{report\_q42.py}
\verbatiminput{../final_version/Q4/part2/code/report_q42.py}

\subsection{问题四-3：波动电价下的多时点调整}

\subsubsection*{q3\_data.py}
\verbatiminput{../final_version/Q4/part3/code/q3_data.py}

\subsubsection*{q3\_forecast.py}
\verbatiminput{../final_version/Q4/part3/code/q3_forecast.py}

\subsubsection*{q3\_model.py}
\verbatiminput{../final_version/Q4/part3/code/q3_model.py}

\subsubsection*{run\_q3.py}
\verbatiminput{../final_version/Q4/part3/code/run_q3.py}

\subsubsection*{verify\_q3.py}
\verbatiminput{../final_version/Q4/part3/code/verify_q3.py}

\subsection{公共模块：第二问基础数据与 LP（q2\_base）}

\subsubsection*{q2\_data.py}
\verbatiminput{../final_version/common/q2_base/q2_data.py}

\subsubsection*{run\_1439.py}
\verbatiminput{../final_version/common/q2_base/run_1439.py}

\subsection{公共模块：附件5结果模板表头（主程序依赖）}

\subsubsection*{template\_layout.py}
\verbatiminput{../final_version/common/efficiency/template_layout.py}

\endgroup